\documentclass[12pt,a4paper]{article}
\usepackage[utf8]{inputenc}
\usepackage[T1]{fontenc}
\usepackage[turkish]{babel}
\usepackage{geometry}
\usepackage{booktabs}
\usepackage{longtable}
\usepackage{array}
\usepackage{hyperref}
\usepackage{xcolor}
\usepackage{titlesec}
\usepackage{enumitem}
\usepackage{fancyhdr}
\usepackage{tabularx}
\usepackage{multirow}
\usepackage{lmodern}
\usepackage{mdframed}
\usepackage{tcolorbox}
\usepackage{graphicx}
\usepackage{amsmath}
\usepackage{colortbl}

\geometry{margin=2.2cm}

% ── Renk paleti ──────────────────────────────────────────────────────────────
\definecolor{critical}{RGB}{192, 0, 0}
\definecolor{high}{RGB}{220, 100, 0}
\definecolor{medium}{RGB}{180, 140, 0}
\definecolor{low}{RGB}{60, 120, 60}
\definecolor{primary}{RGB}{88, 56, 156}
\definecolor{accent}{RGB}{147, 112, 219}
\definecolor{rowgray}{RGB}{248, 248, 250}
\definecolor{headergray}{RGB}{230, 230, 240}

\hypersetup{colorlinks=true, linkcolor=primary, urlcolor=primary}

\titleformat{\section}{\large\bfseries\color{primary}}{}{0em}{}[\titlerule]
\titleformat{\subsection}{\normalsize\bfseries\color{accent}}{}{0em}{\thesubsection\ }

\pagestyle{fancy}
\fancyhf{}
\rhead{\textcolor{gray}{\small CanYoldaşı — Üretim Hazırlık Denetimi}}
\lhead{\textcolor{gray}{\small 16 Temmuz 2026 · GİZLİ}}
\rfoot{\textcolor{gray}{\small \thepage}}
\lfoot{\textcolor{gray}{\small Yalnızca iç kullanım içindir}}

% ── Önem etiketi makroları ──────────────────────────────────────────────────
\newcommand{\CRITICAL}{\colorbox{critical}{\textcolor{white}{\textbf{\small KRİTİK}}}}
\newcommand{\HIGH}{\colorbox{high}{\textcolor{white}{\textbf{\small YÜKSEK}}}}
\newcommand{\MEDIUM}{\colorbox{medium}{\textcolor{white}{\textbf{\small ORTA}}}}
\newcommand{\LOW}{\colorbox{low}{\textcolor{white}{\textbf{\small DÜŞÜK}}}}

\newcommand{\EASY}{\textcolor{low}{\textbf{Kolay}}}
\newcommand{\MEDIUM_D}{\textcolor{medium}{\textbf{Orta}}}
\newcommand{\HARD}{\textcolor{high}{\textbf{Zor}}}

\tcbuselibrary{skins,breakable}
\newtcolorbox{auditbox}[2][]{
  colback=#2!8!white,
  colframe=#2,
  fonttitle=\bfseries\small,
  title=#1,
  breakable,
  left=4pt, right=4pt, top=4pt, bottom=4pt
}

\begin{document}

% ══════════════════════════════════════════════════════════════════════════════
%  KAPAK
% ══════════════════════════════════════════════════════════════════════════════
\begin{center}
  \vspace*{0.5em}
  {\LARGE\bfseries\color{primary} CanYoldaşı}\\[0.4em]
  {\Large\bfseries Üretim Hazırlık Denetimi}\\[0.3em]
  {\normalsize\color{gray} Adım 15 — Derinlemesine Güvenlik ve Kalite Raporu}\\[0.2em]
  {\normalsize\color{gray} 16 Temmuz 2026 · \textbf{GİZLİ — Yalnızca İç Kullanım}}
  \vspace{0.5em}
\end{center}

% ══════════════════════════════════════════════════════════════════════════════
%  YÖNETİCİ ÖZETİ
% ══════════════════════════════════════════════════════════════════════════════
\section{Yönetici Özeti}

\begin{tcolorbox}[colback=primary!6!white, colframe=primary, fonttitle=\bfseries, title=Genel Üretim Hazırlık Skoru]
\begin{center}
  \Huge\bfseries\color{primary} 42 / 100
\end{center}
\vspace{0.3em}
\begin{tabularx}{\textwidth}{lXr}
  Kimlik Doğrulama / Yetkilendirme & & 55/100 \\
  Güvenlik Başlıkları / CORS        & & 20/100 \\
  Hız Sınırlama                     & & 30/100 \\
  Girdi Doğrulama (Zod)             & & 25/100 \\
  Veritabanı Güvenliği / Performans & & 55/100 \\
  Ödeme Güvenliği                   & & 90/100 \\
  Mobil Güvenlik                    & & 45/100 \\
  App Store / Google Play Hazırlığı & & 50/100 \\
  Performans                        & & 50/100 \\
  Hata Yönetimi                     & & 55/100 \\
\end{tabularx}
\end{tcolorbox}

\vspace{0.6em}
\begin{tabularx}{\textwidth}{lX}
  \toprule
  \textbf{Metrik} & \textbf{Tahmin} \\
  \midrule
  App Store'a Kalan Süre & \textbf{4--6 hafta} (kritik engelleyiciler çözülünce) \\
  Google Play'e Kalan Süre & \textbf{4--6 hafta} (App Store ile aynı döngü) \\
  Kritik Engelleyici Sayısı & \textbf{5 adet} \\
  Yüksek Öncelikli Bulgu & 7 adet \\
  Orta Öncelikli Bulgu & 10 adet \\
  Düşük Öncelikli Bulgu & 8 adet \\
  \bottomrule
\end{tabularx}

% ══════════════════════════════════════════════════════════════════════════════
%  KRİTİK ENGELLEYİCİLER
% ══════════════════════════════════════════════════════════════════════════════
\section{Kritik Engelleyiciler (App Store / Google Play Yayınına Engel)}

\begin{auditbox}[C-1 · JWT Secret Sabit Varsayılan Değer]{critical}
\textbf{Problem:} \texttt{JWT\_SECRET} ortam değişkeni ayarlanmamışsa sunucu \texttt{"dev-secret-change-me"} sabit değerini kullanıyor (\texttt{lib/auth.ts}). Üretim ortamında bu değer ayarlanmadan deploy edilirse tüm JWT token'ları tahmin edilebilir ve taklit edilebilir hale gelir.\\[0.3em]
\textbf{Risk:} Herhangi bir kullanıcının hesabı tamamen ele geçirilebilir. Tüm kimlik doğrulama sistemi çöker.\\[0.3em]
\textbf{Önerilen Düzeltme:} Sunucu başlangıcında \texttt{JWT\_SECRET} kontrolü ekle; değer yoksa süreci hata vererek durdur (\textit{fail-fast}). Replit Secrets'ta en az 64 karakter rastgele string olarak tanımla.\\[0.3em]
\textbf{Uygulama Zorluğu:} \EASY\ (30 dakika)
\end{auditbox}

\begin{auditbox}[C-2 · RevenueCat Test Mağaza Anahtarı Üretimde Kullanılıyor]{critical}
\textbf{Problem:} \texttt{artifacts/mobile/services/revenueCat.ts} dosyasında geliştirici uyarısı mevcut: \textit{"EXPO\_PUBLIC\_REVENUECAT\_API\_KEY currently uses the RevenueCat TEST STORE public SDK key. Before release, this MUST be replaced with platform-specific production keys."} Uygulama şu anda test mağaza anahtarıyla çalışıyor; App Store / Google Play geçiş incelemesinde reddedilir.\\[0.3em]
\textbf{Risk:} Gerçek kullanıcı ödemeleri işlenemez. Apple / Google ödeme akışları çalışmaz.\\[0.3em]
\textbf{Önerilen Düzeltme:} RevenueCat dashboard'undan iOS ve Android için ayrı üretim API anahtarları al. \texttt{Platform.select(\{ios: IOS\_KEY, android: ANDROID\_KEY\})} ile platforma özgü başlatma yap. EAS Secrets'a \texttt{EXPO\_PUBLIC\_REVENUECAT\_IOS\_API\_KEY} ve \texttt{EXPO\_PUBLIC\_REVENUECAT\_ANDROID\_API\_KEY} olarak ekle.\\[0.3em]
\textbf{Uygulama Zorluğu:} \EASY\ (1 saat)
\end{auditbox}

\begin{auditbox}[C-3 · Stories Endpoint'inde Kimlik Doğrulama Atlatma]{critical}
\textbf{Problem:} \texttt{POST /api/stories} endpoint'i \texttt{userId} değerini doğrudan \texttt{req.body}'den alıyor ve JWT token'ıyla karşılaştırmıyor. Herhangi bir kimliği doğrulanmış kullanıcı, başka bir kullanıcının \texttt{userId}'siyle hikaye gönderebilir.\\[0.3em]
\textbf{Risk:} Kimliğe bürünme saldırısı. Bir kullanıcı, başka biri adına içerik oluşturabilir.\\[0.3em]
\textbf{Önerilen Düzeltme:} \texttt{stories.ts} route handler'ında \texttt{userId}'yi \texttt{req.body}'den değil \texttt{extractUserId(req)} üzerinden al ve JWT ile zorunlu kıl.\\[0.3em]
\textbf{Uygulama Zorluğu:} \EASY\ (15 dakika)
\end{auditbox}

\begin{auditbox}[C-4 · Güvenlik Başlıkları Eksik (Helmet Yok)]{critical}
\textbf{Problem:} \texttt{app.ts} dosyasında \texttt{helmet} middleware'i yok. XSS koruması, clickjacking önleme (\texttt{X-Frame-Options}), MIME-sniffing önleme (\texttt{X-Content-Type-Options}), HSTS ve CSP başlıkları aktif değil.\\[0.3em]
\textbf{Risk:} Tarayıcı tabanlı saldırılara açık. App Store güvenlik taramasında işaretlenebilir.\\[0.3em]
\textbf{Önerilen Düzeltme:} \texttt{npm install helmet} ile ekle, \texttt{app.use(helmet())} olarak tüm route'lardan önce kayıt et.\\[0.3em]
\textbf{Uygulama Zorluğu:} \EASY\ (20 dakika)
\end{auditbox}

\begin{auditbox}[C-5 · CORS Tüm Origin'lere Açık]{critical}
\textbf{Problem:} \texttt{app.use(cors())} varsayılan yapılandırmasıyla tüm origin'lere (\texttt{*}) izin veriyor. Kötü niyetli bir web sitesi, giriş yapmış kullanıcı adına API istekleri gönderebilir.\\[0.3em]
\textbf{Risk:} CSRF benzeri saldırı vektörü. Yetkisiz cross-origin istekleri.\\[0.3em]
\textbf{Önerilen Düzeltme:} \texttt{cors(\{origin: [process.env.ALLOWED\_ORIGINS]\})} ile beyaz liste tanımla. Mobil app için bu genellikle Expo preview domain'i ve production domain'idir.\\[0.3em]
\textbf{Uygulama Zorluğu:} \EASY\ (30 dakika)
\end{auditbox}

% ══════════════════════════════════════════════════════════════════════════════
%  YÜKSEK ÖNCELİKLİ BULGULAR
% ══════════════════════════════════════════════════════════════════════════════
\section{Yüksek Öncelikli Bulgular}

\begin{auditbox}[H-1 · JWT Token AsyncStorage'da Açık Metin Olarak Tutuluyor]{high}
\textbf{Problem:} \texttt{AuthContext.tsx} JWT token'ı ve kullanıcı profilini \texttt{AsyncStorage}'a (\texttt{@canyoldasi:jwt}) yazıyor. Root'lanmış/jailbreak yapılmış cihazlarda diğer süreçler bu verilere erişebilir.\\[0.3em]
\textbf{Risk:} Cihaz ele geçirilirse tüm kullanıcı oturumları çalınabilir.\\[0.3em]
\textbf{Önerilen Düzeltme:} JWT token'ı için \texttt{expo-secure-store} kullanılmalı. Kullanıcı profil verisi (hassas olmayan) AsyncStorage'da kalabilir.\\[0.3em]
\textbf{Uygulama Zorluğu:} \EASY\ (2 saat)
\end{auditbox}

\begin{auditbox}[H-2 · 401 Yanıtında Otomatik Çıkış Yok]{high}
\textbf{Problem:} Mobil API istemcisinde (\texttt{AuthContext.tsx}) sunucu 401 döndürdüğünde otomatik logout veya token yenileme mekanizması yok. Süresi dolmuş token'lar sessizce hata üretir; kullanıcı anlayamaz.\\[0.3em]
\textbf{Risk:} Kullanıcı deneyimi bozukluğu. Güvenli olmayan oturum durumu.\\[0.3em]
\textbf{Önerilen Düzeltme:} Merkezi \texttt{apiFetch} yardımcısına 401 interceptor ekle; otomatik \texttt{logout()} çağır ve kullanıcıyı login ekranına yönlendir.\\[0.3em]
\textbf{Uygulama Zorluğu:} \EASY\ (1.5 saat)
\end{auditbox}

\begin{auditbox}[H-3 · Beğeni/Yorum Sayaçları Transaction Dışında]{high}
\textbf{Problem:} \texttt{feed.ts} route'unda beğeni ekleme/silme işlemi (INSERT/DELETE \texttt{feed\_likes}) ve \texttt{feed\_posts.likesCount} güncelleme ayrı sorgularda yapılıyor; ortak bir transaction içinde değil. Sunucu arızası arasında oluşursa sayaç kalıcı olarak tutarsız hale gelir.\\[0.3em]
\textbf{Risk:} Veri tutarsızlığı. Yanlış beğeni sayıları.\\[0.3em]
\textbf{Önerilen Düzeltme:} \texttt{db.transaction()} içine al; sayaç artırma/azaltmayı aynı transaction'da yap.\\[0.3em]
\textbf{Uygulama Zorluğu:} \EASY\ (1 saat)
\end{auditbox}

\begin{auditbox}[H-4 · Kritik Sütunlarda Index Eksikliği]{high}
\textbf{Problem:} Sık sorgulanan sütunlarda explicit Drizzle index tanımı yok:\\
\texttt{userId} — feed\_posts, stories, stray\_animals, adoption\_listings, pet\_profiles\\
\texttt{listingId} — adoption\_requests, listing\_promotions, listing\_promotion\_purchases\\
\texttt{status} — stray\_animals, adoption\_listings, adoption\_requests\\
\texttt{createdAt} — hemen tüm tablolar (ORDER BY için kullanılıyor)\\[0.3em]
\textbf{Risk:} Veri büyüdükçe tam tablo taraması; ciddi performans degradasyonu.\\[0.3em]
\textbf{Önerilen Düzeltme:} Drizzle schema'larında \texttt{index("idx\_name").on(t.column)} tanımları ekle; \texttt{executeSql} ile migrate et.\\[0.3em]
\textbf{Uygulama Zorluğu:} \MEDIUM\_D\ (3 saat)
\end{auditbox}

\begin{auditbox}[H-5 · API Route'larında Girdi Doğrulama Neredeyse Yok]{high}
\textbf{Problem:} Zod veya eşdeğer schema doğrulama yalnızca RevenueCat webhook handler'ında (\texttt{rcWebhook.ts}) mevcut. Diğer tüm route'lar \texttt{req.body as Record<string,string>} kalıbıyla çalışıyor; tip güvencesi yok, payload boyutu denetlenmez.\\[0.3em]
\textbf{Risk:} Beklenmedik veri tipleri, aşırı büyük string'ler veya eksik alanlar sunucu hatalarına ya da DB kırılmalarına yol açabilir.\\[0.3em]
\textbf{Önerilen Düzeltme:} \texttt{@workspace/api-zod} aracılığıyla her kritik endpoint için Zod şeması tanımla ve validation middleware olarak ekle. Önce kimlik doğrulama, ilan ve hayvan route'larından başla.\\[0.3em]
\textbf{Uygulama Zorluğu:} \HARD\ (2--3 gün, kapsamlı)
\end{auditbox}

\begin{auditbox}[H-6 · Genel Hız Sınırlama Yok]{high}
\textbf{Problem:} \texttt{POST /api/auth/login}, \texttt{POST /api/auth/register}, \texttt{POST /api/upload} endpoint'leri için hız sınırlama yok (yalnızca şifre sıfırlama kendi custom limiter'ına sahip). Brute-force login ve upload flood saldırılarına açık.\\[0.3em]
\textbf{Risk:} Hesap kırma, servis dışı bırakma (DoS), depolama alanı tükenmesi.\\[0.3em]
\textbf{Önerilen Düzeltme:} \texttt{express-rate-limit} ile en az login (5 deneme/15 dk) ve upload (10 istek/dk) için rate limit ekle.\\[0.3em]
\textbf{Uygulama Zorluğu:} \EASY\ (2 saat)
\end{auditbox}

\begin{auditbox}[H-7 · Android İçin Google Maps API Anahtarı Tanımlanmamış]{high}
\textbf{Problem:} \texttt{app.json} içinde \texttt{android.config.googleMaps.apiKey} alanı yok. Android'de üretim build'inde harita çalışmaz; Google Play incelemesinde tespit edilir.\\[0.3em]
\textbf{Risk:} Uygulamanın temel özelliği (harita) Android'de çalışmaz.\\[0.3em]
\textbf{Önerilen Düzeltme:} Google Cloud Console'da kısıtlanmış bir Android Maps API anahtarı oluştur, \texttt{app.json} \texttt{android.config.googleMaps.apiKey} alanına ekle.\\[0.3em]
\textbf{Uygulama Zorluğu:} \EASY\ (1 saat)
\end{auditbox}

% ══════════════════════════════════════════════════════════════════════════════
%  ORTA ÖNCELİKLİ BULGULAR
% ══════════════════════════════════════════════════════════════════════════════
\section{Orta Öncelikli Bulgular}

\begin{longtable}{p{0.12\textwidth}p{0.24\textwidth}p{0.28\textwidth}p{0.24\textwidth}p{0.07\textwidth}}
\toprule
\textbf{ID} & \textbf{Problem} & \textbf{Risk} & \textbf{Önerilen Düzeltme} & \textbf{Zorluk} \\
\midrule
\endhead

\rowcolor{rowgray}
M-1 &
FlatList büyük feed'lerde \linebreak kullanılıyor &
Hızlı kaydırma sırasında boş alanlar, \linebreak yüksek bellek kullanımı &
\texttt{feed.tsx} ve \texttt{animals.tsx}'te \texttt{@shopify/flash-list} kullan &
\MEDIUM\_D \\

M-2 &
PostCard / AdoptionCard \linebreak memoize edilmemiş &
Liste state değiştiğinde \linebreak tüm kart listesi yeniden render &
\texttt{React.memo()} + \texttt{useCallback} \linebreak renderItem için &
\EASY \\

\rowcolor{rowgray}
M-3 &
Adoption request oluşturma \linebreak transactional değil &
Eşzamanlı isteklerde \linebreak yarış durumu riski &
delete + insert işlemlerini \linebreak \texttt{db.transaction()} içine al &
\EASY \\

M-4 &
Pet şeması \linebreak cascade delete yok &
Pet silinince pet\_posts, \linebreak pet\_health, pet\_followers \linebreak sahipsiz kalır &
\texttt{pets.ts} şemasına \linebreak \texttt{.references(..., \{onDelete:"cascade"\})} ekle &
\EASY \\

\rowcolor{rowgray}
M-5 &
\texttt{uploadImage} fonksiyonu \linebreak 6+ ekranda tekrarlanıyor &
Bakım zorluğu; bir ekranda hata \linebreak diğerlerinde gözden kaçar &
Shared \texttt{utils/uploadImage.ts} \linebreak yardımcısına taşı &
\EASY \\

M-6 &
Üretim console.log'ları \linebreak (\texttt{boost-packages.tsx}, \linebreak \texttt{animal/[id].tsx}) &
Hassas bilgi sızıntısı; \linebreak App Store incelemesinde işaretlenebilir &
\texttt{babel-plugin-transform-remove-console} \linebreak üretim babel config'e ekle &
\EASY \\

\rowcolor{rowgray}
M-7 &
Hardcoded \texttt{loremflickr.com} \linebreak avatar fallback &
Dış CDN bağımlılığı; \linebreak servis kesilirse uygulama bozuk görünür &
Yerel asset (\texttt{assets/images/default-avatar.png}) \linebreak kullan &
\EASY \\

M-8 &
Konum izni reddedilince \linebreak sessiz fallback (İstanbul) &
Yanlış hayvan konumu raporlanabilir; \linebreak kullanıcı bilgilendirilmez &
İzin reddinde açık UI mesajı \linebreak ve "İzin Ver" butonu ekle &
\EASY \\

\rowcolor{rowgray}
M-9 &
Feed discover route \linebreak tüm post'ları JS'e çekiyor &
Veri büyüdükçe bellek/ağ \linebreak performansı bozulur &
Private profil filtrelemesini DB sorgusuna \linebreak (JOIN + WHERE) taşı &
\MEDIUM\_D \\

M-10 &
Sosyal ve feed tablolarında \linebreak \texttt{updated\_at} yok &
Cache invalidation için \linebreak güvenilir timestamp yok &
Eksik tablolara \texttt{updated\_at} \linebreak sütunu ekle + migration &
\EASY \\

\bottomrule
\end{longtable}

% ══════════════════════════════════════════════════════════════════════════════
%  DÜŞÜK ÖNCELİKLİ BULGULAR
% ══════════════════════════════════════════════════════════════════════════════
\section{Düşük Öncelikli Bulgular}

\begin{longtable}{p{0.10\textwidth}p{0.28\textwidth}p{0.25\textwidth}p{0.26\textwidth}p{0.07\textwidth}}
\toprule
\textbf{ID} & \textbf{Problem} & \textbf{Risk} & \textbf{Önerilen Düzeltme} & \textbf{Zorluk} \\
\midrule
\endhead

\rowcolor{rowgray}
L-1 &
\texttt{oauth\_users.email} nullable &
Hesap bağlama tutarsızlığı &
NOT NULL + migration &
\EASY \\

L-2 &
\texttt{social\_profiles.username} \linebreak nullable ama ID olarak kullanılıyor &
Olası null pointer hataları &
NOT NULL kısıtı ekle &
\EASY \\

\rowcolor{rowgray}
L-3 &
Global \texttt{unhandledRejection} \linebreak listener yok &
Yakalanmamış promise hataları \linebreak sessizce geçer &
\texttt{index.ts}'e listener ekle \linebreak ve logla &
\EASY \\

L-4 &
EAS config var, CI/CD yok &
Manuel build riski &
GitHub Actions + EAS build \linebreak pipeline kur &
\HARD \\

\rowcolor{rowgray}
L-5 &
Web stub (\texttt{react-native-maps.web.js}) \linebreak yalnızca geliştirme için gerekli &
Gereksiz dosya; karmaşıklık \linebreak (gereklidir, kaldırmayın) &
— (metro config zorunlu, \linebreak olduğu gibi bırakın) &
— \\

L-6 &
Upload MIME sadece string \linebreak başlangıcını kontrol ediyor &
Magic byte doğrulama yok; \linebreak kötü niyetli dosya yüklenebilir &
\texttt{file-type} paketi ile \linebreak magic byte kontrolü ekle &
\MEDIUM\_D \\

\rowcolor{rowgray}
L-7 &
Bildirim izni için \linebreak explicit akış yok &
iOS'ta sistem iletişim kutusu \linebreak görünmeyebilir &
\texttt{expo-notifications} \linebreak \texttt{requestPermissionsAsync()} akışı kur &
\EASY \\

L-8 &
Icon / Splash screen \linebreak asset boyutları doğrulanmadı &
App Store 1024x1024 ve Android \linebreak adaptive icon gereksinimleri &
Asset boyutlarını kontrol et; \linebreak App Store Connect'e yükle ve denetle &
\EASY \\

\bottomrule
\end{longtable}

% ══════════════════════════════════════════════════════════════════════════════
%  KONU BAZLI DETAYLAR
% ══════════════════════════════════════════════════════════════════════════════
\section{Konu Bazlı Detaylı Değerlendirme}

\subsection*{Ödeme Güvenliği (RevenueCat + Stripe)}

\begin{tcolorbox}[colback=low!8!white, colframe=low, fonttitle=\bfseries, title=İyi Uygulama — Güçlü Yönler]
\begin{itemize}[leftmargin=1em]
  \item Atomik transaction ile verify-purchase pipeline (yarış durumu korumalı)
  \item Sunucu tarafı \texttt{PACKAGE\_MAP} — istemci \texttt{durationDays} değerini taklit edemez
  \item \texttt{transaction\_identifier} UNIQUE index — çift harcama önlenir
  \item RevenueCat webhook: iki katmanlı güvenlik (Auth token + HMAC-SHA256)
  \item 5 dakikalık imza geçerlilik penceresi
  \item Idempotent webhook işleme (\texttt{revenuecat\_event\_id} UNIQUE)
  \item Ledger tabanlı promosyon hesaplama (stacking-safe refund)
\end{itemize}
\end{tcolorbox}

\subsection*{Veritabanı Performans Özeti}

Kritik tablolar ve önerilen index'ler:

\begin{tabularx}{\textwidth}{lXl}
\toprule
\textbf{Tablo} & \textbf{Eksik Index'ler} & \textbf{Öncelik} \\
\midrule
\texttt{feed\_posts} & \texttt{userId}, \texttt{createdAt}, \texttt{status} & Yüksek \\
\texttt{stray\_animals} & \texttt{userId}, \texttt{status}, \texttt{createdAt} & Yüksek \\
\texttt{adoption\_listings} & \texttt{userId}, \texttt{status}, \texttt{promotedUntil} & Yüksek \\
\texttt{listing\_promotion\_purchases} & \texttt{listingId}, \texttt{userId}, \texttt{purchaseStatus} & Yüksek \\
\texttt{adoption\_requests} & \texttt{listingId}, \texttt{requesterId}, \texttt{status} & Yüksek \\
\texttt{pet\_profiles} & \texttt{userId}, \texttt{createdAt} & Orta \\
\texttt{stories} & \texttt{userId}, \texttt{createdAt} & Orta \\
\bottomrule
\end{tabularx}

\subsection*{Mobil Güvenlik Özeti}

\begin{tabularx}{\textwidth}{lXl}
\toprule
\textbf{Alan} & \textbf{Durum} & \textbf{Öncelik} \\
\midrule
JWT depolama & AsyncStorage (şifresiz) & Yüksek \\
401 interceptor & Yok & Yüksek \\
Deep link koruması & Kısmi & Orta \\
Form validasyonu & Tutarsız & Orta \\
XSS sanitizasyon & Yok (React Native'de risk düşük) & Düşük \\
RevenueCat üretim anahtarı & Test anahtarı kullanılıyor & \textbf{Kritik} \\
console.log'lar & Kısmen \texttt{\_\_DEV\_\_} korumalı & Orta \\
\bottomrule
\end{tabularx}

% ══════════════════════════════════════════════════════════════════════════════
%  APP STORE / GOOGLE PLAY HAZIRLIK LİSTESİ
% ══════════════════════════════════════════════════════════════════════════════
\section{App Store / Google Play Hazırlık Kontrol Listesi}

\begin{longtable}{p{0.55\textwidth}p{0.15\textwidth}p{0.25\textwidth}}
\toprule
\textbf{Gereksinim} & \textbf{Durum} & \textbf{Not} \\
\midrule
\endhead
Bundle ID tanımlanmış & \textcolor{low}{\textbf{TAMAM}} & \texttt{com.canyoldasi.app} \\
\rowcolor{rowgray}
Versiyon numarası ayarlı & \textcolor{low}{\textbf{TAMAM}} & 1.0.0 \\
iOS build numarası & \textcolor{low}{\textbf{TAMAM}} & 1.0.0.1 \\
\rowcolor{rowgray}
Android versionCode & \textcolor{low}{\textbf{TAMAM}} & 1 \\
iOS izin açıklamaları (Info.plist) & \textcolor{low}{\textbf{TAMAM}} & Konum, kamera, galeri \\
\rowcolor{rowgray}
Android izin bildirimleri & \textcolor{low}{\textbf{TAMAM}} & Konum, kamera \\
Privacy Manifest (iOS 17+) & \textcolor{low}{\textbf{TAMAM}} & NSUserTrackingUsageDescription mevcut \\
\rowcolor{rowgray}
ITSAppUsesNonExemptEncryption & \textcolor{low}{\textbf{TAMAM}} & false \\
EAS build konfigürasyonu & \textcolor{low}{\textbf{TAMAM}} & eas.json mevcut \\
\rowcolor{rowgray}
RevenueCat üretim anahtarları & \textcolor{critical}{\textbf{EKSİK}} & \textbf{KRİTİK ENGELLEYICI} \\
Google Maps API anahtarı (Android) & \textcolor{critical}{\textbf{EKSİK}} & \textbf{KRİTİK ENGELLEYICI} \\
\rowcolor{rowgray}
JWT Secret üretim değeri & \textcolor{critical}{\textbf{RİSKLİ}} & Varsayılan değer tehlikeli \\
App Store 1024x1024 icon & \textcolor{medium}{\textbf{DOĞRULA}} & Boyut kontrol edilmedi \\
\rowcolor{rowgray}
Android adaptive icon & \textcolor{medium}{\textbf{DOĞRULA}} & Boyut kontrol edilmedi \\
Splash screen & \textcolor{low}{\textbf{TAMAM}} & app.json'da tanımlı \\
\rowcolor{rowgray}
CI/CD pipeline & \textcolor{medium}{\textbf{EKSİK}} & Manuel build riski var \\
console.log production strip & \textcolor{medium}{\textbf{EKSİK}} & Babel plugin gerekli \\
\rowcolor{rowgray}
Gizlilik politikası URL'si & \textcolor{medium}{\textbf{EKSİK}} & Her iki mağaza zorunlu kılıyor \\
Kullanım koşulları URL'si & \textcolor{medium}{\textbf{EKSİK}} & Her iki mağaza zorunlu kılıyor \\
\rowcolor{rowgray}
Yaş derecelendirmesi & \textcolor{medium}{\textbf{DOĞRULA}} & App Store Connect'te ayarlanmalı \\
\bottomrule
\end{longtable}

% ══════════════════════════════════════════════════════════════════════════════
%  ÖNERİLEN ÇALIŞMA SIRASI
% ══════════════════════════════════════════════════════════════════════════════
\section{Önerilen Çalışma Sırası (Sprint Planı)}

\begin{longtable}{p{0.08\textwidth}p{0.15\textwidth}p{0.52\textwidth}p{0.12\textwidth}}
\toprule
\textbf{Hafta} & \textbf{Kapsam} & \textbf{Görevler} & \textbf{Etki} \\
\midrule
\endhead

\rowcolor{critical!10}
1 &
Kritik \linebreak Güvenlik &
C-1: JWT fail-fast · C-2: RC üretim anahtarları · C-3: Stories auth bypass · C-4: Helmet · C-5: CORS kısıtla &
\textcolor{critical}{\textbf{Kritik}} \\

\rowcolor{high!10}
2 &
Yüksek \linebreak Güvenlik &
H-1: expo-secure-store · H-2: 401 interceptor · H-3: Transaction'lar · H-6: Rate limiting · H-7: Google Maps key &
\textcolor{high}{\textbf{Yüksek}} \\

\rowcolor{medium!10}
3 &
DB \& \linebreak Performans &
H-4: Index'ler · M-9: Feed sorgu optimizasyonu · M-10: updated\_at · M-3/M-4: Transaction/cascade &
\textcolor{medium}{\textbf{Orta}} \\

\rowcolor{medium!10}
4 &
Mobil \linebreak Kalite &
M-1: FlashList · M-2: Memoization · M-5: uploadImage utils · M-6: console.log · M-7: local avatar · M-8: izin UI &
\textcolor{medium}{\textbf{Orta}} \\

\rowcolor{low!10}
5--6 &
Validasyon \linebreak \& Store &
H-5: Zod validation tüm route'lar · Gizlilik politikası · Kullanım koşulları · App Store Connect metadata · EAS CI/CD &
\textcolor{low}{\textbf{Düşük}} \\

\bottomrule
\end{longtable}

% ══════════════════════════════════════════════════════════════════════════════
%  SONUÇ
% ══════════════════════════════════════════════════════════════════════════════
\section{Sonuç}

CanYoldaşı, ödeme güvenliği açısından (RevenueCat + Stripe entegrasyonu) endüstri standardında bir uygulama sunmaktadır. Atomik transaction'lar, HMAC doğrulama ve idempotent webhook işleme bu alanlarda güçlü bir temel oluşturmaktadır.

Bununla birlikte, \textbf{5 kritik ve 7 yüksek öncelikli bulgu} App Store / Google Play yayını öncesinde giderilmesi zorunlu sorunlar içermektedir. En acil müdahale gerektiren konular şunlardır:

\begin{enumerate}
  \item \textbf{RevenueCat test anahtarının üretim anahtarıyla değiştirilmesi} — ödeme akışları çalışmaz.
  \item \textbf{JWT Secret varsayılan değerinin kaldırılması} — tüm kimlik doğrulama sistemi açık.
  \item \textbf{Helmet ve CORS konfigürasyonu} — temel web güvenliği eksik.
  \item \textbf{Android Google Maps API anahtarı} — temel özellik çalışmaz.
\end{enumerate}

Bu engeller giderildikten sonra Hafta 5--6'daki validasyon ve mağaza metadata çalışmalarıyla uygulama \textbf{4--6 hafta içinde} App Store ve Google Play'de yayınlanmaya hazır hale gelebilir.

\vspace{1.5em}
\noindent\rule{\textwidth}{0.4pt}
\begin{center}
  \small\color{gray}
  CanYoldaşı Üretim Hazırlık Denetimi · Adım 15 · 16 Temmuz 2026\\
  pnpm monorepo · Expo SDK 54 · Express 5 · PostgreSQL · RevenueCat · Stripe\\
  \textit{Bu belge otomatik analiz araçlarıyla üretilmiştir. Tüm bulgular geliştirici incelemesine tabidir.}
\end{center}

\end{document}
